\section{Quotients of triangulated categories}
\label{section-quotients}

\noindent
Given a triangulated category and a triangulated subcategory we can
construct another triangulated category by taking the ``quotient''.
The construction uses a localization. This is similar to the quotient
of an abelian category by a Serre subcategory, see
Homology, Section \ref{homology-section-serre-subcategories}.
Before we do the actual construction we briefly discuss kernels
of exact functors.

\begin{definition}
\label{definition-saturated}
Let $\mathcal{D}$ be a pre-triangulated category. We say a full
pre-triangulated subcategory $\mathcal{D}'$ of $\mathcal{D}$ is
{\it saturated} if whenever $X \oplus Y$ is isomorphic to an object
of $\mathcal{D}'$ then both $X$ and $Y$ are isomorphic to objects
of $\mathcal{D}'$.
\end{definition}

\noindent
A saturated triangulated subcategory is sometimes called a
{\it thick triangulated subcategory}. In some references, this is
only used for strictly full triangulated subcategories (and sometimes
the definition is written such that it implies strictness).
There is another notion, that of an {\it \'epaisse triangulated
subcategory}. The definition is that given a commutative diagram
$$
\xymatrix{
& S \ar[rd] \\
X \ar[ru] \ar[rr] & & Y \ar[r] & T \ar[r] & X[1]
}
$$
where the second line is a distinguished triangle and $S$ and $T$ are isomorphic
to objects of $\mathcal{D}'$, then also $X$ and $Y$ are isomorphic to objects
of $\mathcal{D}'$. It turns out that this is equivalent to being
saturated (this is elementary and can be found in \cite{Rickard-derived})
and the notion of a saturated category is easier to work with.

\begin{lemma}
\label{lemma-triangle-functor-kernel}
Let $F : \mathcal{D} \to \mathcal{D}'$ be an exact functor of
pre-triangulated categories. Let $\mathcal{D}''$ be the full subcategory
of $\mathcal{D}$ with objects
$$
\Ob(\mathcal{D}'') =
\{X \in \Ob(\mathcal{D}) \mid F(X) = 0\}
$$
Then $\mathcal{D}''$ is a strictly full saturated pre-triangulated
subcategory of $\mathcal{D}$. If $\mathcal{D}$ is a triangulated category,
then $\mathcal{D}''$ is a triangulated subcategory.
\end{lemma}

\begin{proof}
It is clear that $\mathcal{D}''$ is preserved under $[1]$ and $[-1]$.
If $(X, Y, Z, f, g, h)$ is a distinguished triangle of $\mathcal{D}$
and $F(X) = F(Y) = 0$, then also $F(Z) = 0$ as
$(F(X), F(Y), F(Z), F(f), F(g), F(h))$ is distinguished.
Hence we may apply
Lemma \ref{lemma-triangulated-subcategory}
to see that $\mathcal{D}''$ is a pre-triangulated subcategory (respectively
a triangulated subcategory if $\mathcal{D}$ is a triangulated category).
The final assertion of being saturated follows from
$F(X) \oplus F(Y) = 0 \Rightarrow F(X) = F(Y) = 0$.
\end{proof}

\begin{lemma}
\label{lemma-homological-functor-kernel}
Let $H : \mathcal{D} \to \mathcal{A}$ be a homological functor of
a pre-triangulated category into an abelian category.
Let $\mathcal{D}'$ be the full subcategory of $\mathcal{D}$ with objects
$$
\Ob(\mathcal{D}') =
\{X \in \Ob(\mathcal{D}) \mid
H(X[n]) = 0\text{ for all }n \in \mathbf{Z}\}
$$
Then $\mathcal{D}'$ is a strictly full saturated pre-triangulated subcategory
of $\mathcal{D}$. If $\mathcal{D}$ is a triangulated category, then
$\mathcal{D}'$ is a triangulated subcategory.
\end{lemma}

\begin{proof}
It is clear that $\mathcal{D}'$ is preserved under $[1]$ and $[-1]$.
If $(X, Y, Z, f, g, h)$ is a distinguished triangle of $\mathcal{D}$
and $H(X[n]) = H(Y[n]) = 0$ for all $n$, then also $H(Z[n]) = 0$ for all $n$
by the long exact sequence (\ref{equation-long-exact-cohomology-sequence}).
Hence we may apply
Lemma \ref{lemma-triangulated-subcategory}
to see that $\mathcal{D}'$ is a pre-triangulated subcategory (respectively
a triangulated subcategory if $\mathcal{D}$ is a triangulated category).
The assertion of being saturated follows from
\begin{align*}
H((X \oplus Y)[n]) = 0 & \Rightarrow H(X[n] \oplus Y[n]) = 0 \\
& \Rightarrow H(X[n]) \oplus H(Y[n]) = 0 \\
& \Rightarrow H(X[n]) = H(Y[n]) = 0
\end{align*}
for all $n \in \mathbf{Z}$.
\end{proof}

\begin{lemma}
\label{lemma-homological-functor-bounded}
Let $H : \mathcal{D} \to \mathcal{A}$ be a homological functor of
a pre-triangulated category into an abelian category.
Let $\mathcal{D}_H^{+}, \mathcal{D}_H^{-}, \mathcal{D}_H^b$
be the full subcategory of $\mathcal{D}$ with objects
$$
\begin{matrix}
\Ob(\mathcal{D}_H^{+}) =
\{X \in \Ob(\mathcal{D}) \mid
H(X[n]) = 0\text{ for all }n \ll 0\} \\
\Ob(\mathcal{D}_H^{-}) =
\{X \in \Ob(\mathcal{D}) \mid
H(X[n]) = 0\text{ for all }n \gg 0\} \\
\Ob(\mathcal{D}_H^b) =
\{X \in \Ob(\mathcal{D}) \mid
H(X[n]) = 0\text{ for all }|n| \gg 0\}
\end{matrix}
$$
Each of these is a strictly full saturated pre-triangulated subcategory
of $\mathcal{D}$. If $\mathcal{D}$ is a triangulated category, then
each is a triangulated subcategory.
\end{lemma}

\begin{proof}
Let us prove this for $\mathcal{D}_H^{+}$.
It is clear that it is preserved under $[1]$ and $[-1]$.
If $(X, Y, Z, f, g, h)$ is a distinguished triangle of $\mathcal{D}$
and $H(X[n]) = H(Y[n]) = 0$ for all $n \ll 0$, then also $H(Z[n]) = 0$
for all $n \ll 0$ by the long exact sequence
(\ref{equation-long-exact-cohomology-sequence}).
Hence we may apply
Lemma \ref{lemma-triangulated-subcategory}
to see that $\mathcal{D}_H^{+}$ is a pre-triangulated subcategory
(respectively a triangulated subcategory if $\mathcal{D}$ is a
triangulated category). The assertion of being saturated follows from
\begin{align*}
H((X \oplus Y)[n]) = 0 & \Rightarrow H(X[n] \oplus Y[n]) = 0 \\
& \Rightarrow H(X[n]) \oplus H(Y[n]) = 0 \\
& \Rightarrow H(X[n]) = H(Y[n]) = 0
\end{align*}
for all $n \in \mathbf{Z}$.
\end{proof}

\begin{definition}
\label{definition-kernel-category}
Let $\mathcal{D}$ be a (pre-)triangulated category.
\begin{enumerate}
\item Let $F : \mathcal{D} \to \mathcal{D}'$ be an exact functor.
The {\it kernel of $F$} is the strictly full saturated
(pre-)triangulated subcategory described in
Lemma \ref{lemma-triangle-functor-kernel}.
\item Let $H : \mathcal{D} \to \mathcal{A}$ be a homological functor.
The {\it kernel of $H$} is the strictly full saturated
(pre-)triangulated subcategory described in
Lemma \ref{lemma-homological-functor-kernel}.
\end{enumerate}
These are sometimes denoted $\Ker(F)$ or $\Ker(H)$.
\end{definition}

\noindent
The proof of the following lemma uses TR4.

\begin{lemma}
\label{lemma-construct-multiplicative-system}
Let $\mathcal{D}$ be a triangulated category.
Let $\mathcal{D}' \subset \mathcal{D}$ be a full triangulated
subcategory. Set
\begin{equation}
\label{equation-multiplicative-system}
S =
\left\{
\begin{matrix}
f \in \text{Arrows}(\mathcal{D})
\text{ such that there exists a distinguished triangle }\\
(X, Y, Z, f, g, h) \text{ of }\mathcal{D}\text{ with }
Z\text{ isomorphic to an object of }\mathcal{D}'
\end{matrix}
\right\}
\end{equation}
Then $S$ is a multiplicative system compatible with the triangulated
structure on $\mathcal{D}$. In this situation the following are equivalent
\begin{enumerate}
\item $S$ is a saturated multiplicative system,
\item $\mathcal{D}'$ is a saturated triangulated subcategory.
\end{enumerate}
\end{lemma}

\begin{proof}
To prove the first assertion we have to prove that
MS1, MS2, MS3 and MS5, MS6 hold.

\medskip\noindent
Proof of MS1. It is clear that identities are in $S$ because
$(X, X, 0, 1, 0, 0)$ is distinguished for every object $X$ of $\mathcal{D}$
and because $0$ is an object of $\mathcal{D}'$. Let $f : X \to Y$
and $g : Y \to Z$ be composable morphisms contained in $S$.
Choose distinguished triangles $(X, Y, Q_1, f, p_1, d_1)$,
$(X, Z, Q_2, g \circ f, p_2, d_2)$, and $(Y, Z, Q_3, g, p_3, d_3)$.
By assumption we know that $Q_1$ and $Q_3$ are isomorphic to objects
of $\mathcal{D}'$. By TR4 we know there exists a distinguished
triangle $(Q_1, Q_2, Q_3, a, b, c)$. Since $\mathcal{D}'$ is a
triangulated subcategory we conclude that $Q_2$ is isomorphic to
an object of $\mathcal{D}'$. Hence $g \circ f \in S$.

\medskip\noindent
Proof of MS3. Let $a : X \to Y$ be a morphism and let $t : Z \to X$ be
an element of $S$ such that $a \circ t = 0$. To prove LMS3 it suffices to
find an $s \in S$ such that $s \circ a = 0$, compare with the proof of
Lemma \ref{lemma-triangle-functor-localize}. Choose a distinguished
triangle $(Z, X, Q, t, g, h)$ using TR1 and TR2. Since $a \circ t = 0$
we see by
Lemma \ref{lemma-representable-homological}
there exists a morphism $i : Q \to Y$ such that $i \circ g = a$.
Finally, using TR1 again we can choose a triangle
$(Q, Y, W, i, s, k)$. Here is a picture
$$
\xymatrix{
Z \ar[r]_t & X \ar[r]_g \ar[d]^1 & Q \ar[r] \ar[d]^i & Z[1] \\
& X \ar[r]_a & Y \ar[d]^s \\
& & W
}
$$
Since $t \in S$ we see that $Q$ is isomorphic to an object of $\mathcal{D}'$.
Hence $s \in S$. Finally, $s \circ a = s \circ i \circ g = 0$ as
$s \circ i = 0$ by Lemma \ref{lemma-composition-zero}.
We conclude that LMS3 holds.
The proof of RMS3 is dual.

\medskip\noindent
Proof of MS5. Follows as distinguished triangles and $\mathcal{D}'$
are stable under translations

\medskip\noindent
Proof of MS6. Suppose given a commutative diagram
$$
\xymatrix{
X \ar[r] \ar[d]^s &
Y \ar[d]^{s'} \\
X' \ar[r] &
Y'
}
$$
with $s, s' \in S$. By
Proposition \ref{proposition-9}
we can extend this to a nine square diagram. As $s, s'$ are elements of $S$
we see that $X'', Y''$ are isomorphic to objects of $\mathcal{D}'$.
Since $\mathcal{D}'$ is a full triangulated subcategory we see that
$Z''$ is also isomorphic to an object of $\mathcal{D}'$.
Whence the morphism $Z \to Z'$
is an element of $S$. This proves MS6.

\medskip\noindent
MS2 is a formal consequence of MS1, MS5, and MS6, see
Lemma \ref{lemma-localization-conditions}.
This finishes the proof of the first assertion of the lemma.

\medskip\noindent
Let's assume that $S$ is saturated. (In the following we will use
rotation of distinguished triangles without further mention.)
Let $X \oplus Y$ be an object isomorphic to an object of $\mathcal{D}'$.
Consider the morphism $f : 0 \to X$. The composition
$0 \to X \to X \oplus Y$ is an element
of $S$ as $(0, X \oplus Y, X \oplus Y, 0, 1, 0)$ is a distinguished
triangle. The composition $Y[-1] \to 0 \to X$ is an element of $S$
as $(X, X \oplus Y, Y, (1, 0), (0, 1), 0)$ is a distinguished triangle, see
Lemma \ref{lemma-split}.
Hence $0 \to X$ is an element of $S$ (as $S$ is saturated).
Thus $X$ is isomorphic to an object of $\mathcal{D}'$ as desired.

\medskip\noindent
Finally, assume $\mathcal{D}'$ is a saturated triangulated subcategory.
Let
$$
W \xrightarrow{h}
X \xrightarrow{g}
Y \xrightarrow{f} Z
$$
be composable morphisms of $\mathcal{D}$ such that $fg, gh \in S$.
We will build up a picture of objects as in the diagram below.
$$
\xymatrix{
 & &
Q_{12} \ar[rd] & &
Q_{23} \ar[rd] \\
 &
Q_1 \ar[ld]_{\! + \! 1} \ar[ru] & &
Q_2 \ar[ld]_{\! + \! 1} \ar[ll]_{\! + \! 1} \ar[ru] & &
Q_3 \ar[ld]_{\! + \! 1} \ar[ll]_{\! + \! 1} \\
W \ar[rr] & &
X \ar[lu] \ar[rr] & &
Y \ar[lu] \ar[rr] & &
Z \ar[lu]
}
$$
First choose distinguished triangles
$(W, X, Q_1)$, $(X, Y, Q_2)$, $(Y, Z, Q_3)$ $(W, Y, Q_{12})$, and
$(X, Z, Q_{23})$. Denote $s : Q_2 \to Q_1[1]$ the composition
$Q_2 \to X[1] \to Q_1[1]$. Denote $t : Q_3 \to Q_2[1]$ the
composition $Q_3 \to Y[1] \to Q_2[1]$.
By TR4 applied to the composition $W \to X \to Y$
and the composition $X \to Y \to Z$ there exist
distinguished triangles $(Q_1, Q_{12}, Q_2)$ and $(Q_2, Q_{23}, Q_3)$
which use the morphisms $s$ and $t$.
The objects $Q_{12}$ and $Q_{23}$ are isomorphic to objects of
$\mathcal{D}'$ as $W \to Y$ and $X \to Z$ are assumed in $S$.
Hence also $s[1]t$ is an element of $S$ as $S$ is closed under compositions
and shifts.
Note that $s[1]t = 0$ as $Y[1] \to Q_2[1] \to X[2]$ is zero, see
Lemma \ref{lemma-composition-zero}.
Hence $Q_3[1] \oplus Q_1[2]$ is isomorphic to an object of
$\mathcal{D}'$, see Lemma \ref{lemma-split}.
By assumption on $\mathcal{D}'$ we conclude that $Q_3$ and $Q_1$ are isomorphic
to objects of $\mathcal{D}'$. Looking at the distinguished triangle
$(Q_1, Q_{12}, Q_2)$ we conclude that $Q_2$ is also isomorphic to
an object of $\mathcal{D}'$. Looking at the distinguished triangle
$(X, Y, Q_2)$ we finally conclude that $g \in S$. (It is also
follows that $h, f \in S$, but we don't need this.)
\end{proof}

\begin{definition}
\label{definition-quotient-category}
Let $\mathcal{D}$ be a triangulated category.
Let $\mathcal{B}$ be a full triangulated subcategory.
We define the {\it quotient category $\mathcal{D}/\mathcal{B}$}
by the formula $\mathcal{D}/\mathcal{B} = S^{-1}\mathcal{D}$, where
$S$ is the multiplicative system of $\mathcal{D}$ associated to
$\mathcal{B}$ via
Lemma \ref{lemma-construct-multiplicative-system}.
The localization functor $Q : \mathcal{D} \to \mathcal{D}/\mathcal{B}$
is called the {\it quotient functor} in this case.
\end{definition}

\noindent
Note that the quotient functor
$Q : \mathcal{D} \to \mathcal{D}/\mathcal{B}$
is an exact functor of triangulated categories, see
Proposition \ref{proposition-construct-localization}.
The universal property of this construction is the following.

\begin{lemma}
\label{lemma-universal-property-quotient}
\begin{slogan}
The universal property of the Verdier quotient.
\end{slogan}
Let $\mathcal{D}$ be a triangulated category. Let $\mathcal{B}$
be a full triangulated subcategory of $\mathcal{D}$. Let
$Q : \mathcal{D} \to \mathcal{D}/\mathcal{B}$ be the quotient functor.
\begin{enumerate}
\item If $H : \mathcal{D} \to \mathcal{A}$ is a homological functor into
an abelian category $\mathcal{A}$ such that
$\mathcal{B} \subset \Ker(H)$ then there exists a unique factorization
$H' : \mathcal{D}/\mathcal{B} \to \mathcal{A}$ such that $H = H' \circ Q$
and $H'$ is a homological functor too.
\item If $F : \mathcal{D} \to \mathcal{D}'$ is an exact functor into
a pre-triangulated category $\mathcal{D}'$ such that
$\mathcal{B} \subset \Ker(F)$ then there exists a unique factorization
$F' : \mathcal{D}/\mathcal{B} \to \mathcal{D}'$ such that $F = F' \circ Q$
and $F'$ is an exact functor too.
\end{enumerate}
\end{lemma}

\begin{proof}
This lemma follows from
Lemma \ref{lemma-universal-property-localization}.
Namely, if $f : X \to Y$ is a morphism of $\mathcal{D}$
such that for some distinguished triangle $(X, Y, Z, f, g, h)$
the object $Z$ is isomorphic to an object of $\mathcal{B}$, then
$H(f)$, resp.\ $F(f)$ is an isomorphism under the assumptions of
(1), resp.\ (2). Details omitted.
\end{proof}

\noindent
The kernel of the quotient functor can be described as follows.

\begin{lemma}
\label{lemma-kernel-quotient}
Let $\mathcal{D}$ be a triangulated category.
Let $\mathcal{B}$ be a full triangulated subcategory.
The kernel of the quotient functor
$Q : \mathcal{D} \to \mathcal{D}/\mathcal{B}$
is the strictly full subcategory of $\mathcal{D}$ whose objects are
$$
\Ob(\Ker(Q)) =
\left\{
\begin{matrix}
Z \in \Ob(\mathcal{D})
\text{ such that there exists a }Z' \in \Ob(\mathcal{D}) \\
\text{ such that }Z \oplus Z'\text{ is isomorphic to an object of }\mathcal{B}
\end{matrix}
\right\}
$$
In other words it is the smallest strictly full saturated triangulated
subcategory of $\mathcal{D}$ containing $\mathcal{B}$.
\end{lemma}

\begin{proof}
First note that the kernel is automatically a strictly full
triangulated subcategory containing summands of any of its objects, see
Lemma \ref{lemma-triangle-functor-kernel}.
The description of its objects follows from the definitions and
Lemma \ref{lemma-kernel-localization} part (4).
\end{proof}

\noindent
Let $\mathcal{D}$ be a triangulated category.
At this point we have constructions which induce order
preserving maps between
\begin{enumerate}
\item the partially ordered set of multiplicative systems $S$ in $\mathcal{D}$
compatible with the triangulated structure, and
\item the partially ordered set of full triangulated subcategories
$\mathcal{B} \subset \mathcal{D}$.
\end{enumerate}
Namely, the constructions are given by
$S \mapsto \mathcal{B}(S) = \Ker(Q : \mathcal{D} \to S^{-1}\mathcal{D})$
and $\mathcal{B} \mapsto S(\mathcal{B})$
where $S(\mathcal{B})$ is the multiplicative set of
(\ref{equation-multiplicative-system}), i.e.,
$$
S(\mathcal{B}) =
\left\{
\begin{matrix}
f \in \text{Arrows}(\mathcal{D})
\text{ such that there exists a distinguished triangle }\\
(X, Y, Z, f, g, h) \text{ of }\mathcal{D}\text{ with }
Z\text{ isomorphic to an object of }\mathcal{B}
\end{matrix}
\right\}
$$
Note that it is not the case that these operations are mutually inverse.

\begin{lemma}
\label{lemma-operations}
Let $\mathcal{D}$ be a triangulated category. The operations described above
have the following properties
\begin{enumerate}
\item $S(\mathcal{B}(S))$ is the ``saturation'' of $S$, i.e., it is the
smallest saturated multiplicative system in $\mathcal{D}$ containing $S$, and
\item $\mathcal{B}(S(\mathcal{B}))$ is the ``saturation'' of $\mathcal{B}$,
i.e., it is the smallest strictly full saturated triangulated subcategory of
$\mathcal{D}$ containing $\mathcal{B}$.
\end{enumerate}
In particular, the constructions define mutually inverse maps between
the (partially ordered) set of saturated multiplicative systems in
$\mathcal{D}$ compatible with the triangulated structure on $\mathcal{D}$
and
the (partially ordered) set of strictly full saturated triangulated
subcategories of $\mathcal{D}$.
\end{lemma}

\begin{proof}
First, let's start with a full triangulated subcategory $\mathcal{B}$. Then
$\mathcal{B}(S(\mathcal{B})) =
\Ker(Q : \mathcal{D} \to \mathcal{D}/\mathcal{B})$
and hence (2) is the content of
Lemma \ref{lemma-kernel-quotient}.

\medskip\noindent
Next, suppose that $S$ is multiplicative system in $\mathcal{D}$ compatible
with the triangulation on $\mathcal{D}$. Then
$\mathcal{B}(S) = \Ker(Q : \mathcal{D} \to S^{-1}\mathcal{D})$.
Hence (using
Lemma \ref{lemma-third-object-zero}
in the localized category)
\begin{align*}
S(\mathcal{B}(S))
& =
\left\{
\begin{matrix}
f \in \text{Arrows}(\mathcal{D})
\text{ such that there exists a distinguished}\\
\text{triangle }(X, Y, Z, f, g, h) \text{ of }\mathcal{D}\text{ with }Q(Z) = 0
\end{matrix}
\right\}.
\\
& =
\{f \in \text{Arrows}(\mathcal{D}) \mid Q(f)\text{ is an isomorphism}\} \\
& = \hat S = S'
\end{align*}
in the notation of
Categories, Lemma \ref{categories-lemma-what-gets-inverted}.
The final statement of that lemma finishes the proof.
\end{proof}

\begin{lemma}
\label{lemma-acyclic-general}
Let $H : \mathcal{D} \to \mathcal{A}$ be a homological functor from a
triangulated category $\mathcal{D}$ to an abelian category $\mathcal{A}$, see
Definition \ref{definition-homological}.
The subcategory $\Ker(H)$ of $\mathcal{D}$ is a strictly full
saturated triangulated subcategory of $\mathcal{D}$ whose corresponding
saturated multiplicative system (see
Lemma \ref{lemma-operations})
is the set
$$
S = \{f \in \text{Arrows}(\mathcal{D}) \mid
H^i(f)\text{ is an isomorphism for all }i \in \mathbf{Z}\}.
$$
The functor $H$ factors through the quotient functor
$Q : \mathcal{D} \to \mathcal{D}/\Ker(H)$.
\end{lemma}

\begin{proof}
The category $\Ker(H)$ is a strictly full saturated triangulated
subcategory of $\mathcal{D}$ by
Lemma \ref{lemma-homological-functor-kernel}.
The set $S$ is a saturated multiplicative system compatible with the
triangulated structure by
Lemma \ref{lemma-homological-functor-localize}.
Recall that the multiplicative system corresponding to
$\Ker(H)$ is the set
$$
\left\{
\begin{matrix}
f \in \text{Arrows}(\mathcal{D})
\text{ such that there exists a distinguished triangle }\\
(X, Y, Z, f, g, h)\text{ with } H^i(Z) = 0 \text{ for all }i
\end{matrix}
\right\}.
$$
By the long exact cohomology sequence, see
(\ref{equation-long-exact-cohomology-sequence}),
it is clear that $f$ is an element of this set if and only if $f$ is
an element of $S$. Finally, the factorization of $H$ through $Q$ is a
consequence of
Lemma \ref{lemma-universal-property-quotient}.
\end{proof}










